\hypertarget{appendix-18b8dd966b5d}{}
\section*{Bfloat16}
\section{Bfloat16}\label{bfloat16}

\textbf{Curriculum depth:} Developed · \textbf{First seen:} C006 · \textbf{Latest development:} C006

\subsection{Definition and mechanism}\label{definition-and-mechanism}

BF16 has1/8/7 fields, bias127 and8 normal precision bits. It keeps broad exponent range at coarse spacing2\^{}-7 above1. A nearest-even conversion must handle NaNs separately from finite rounding.

\subsection{Prerequisites}\label{prerequisites}

\begin{itemize}
\tightlist
\item
  \hyperref[appendix-343615f49ddf]{IEEE-754 floating point} --- supplies the representation or validity rule used above.
\item
  \hyperref[appendix-614afe1de592]{Rounding to nearest ties to even} --- supplies the representation or validity rule used above.
\end{itemize}

\subsection{Used by}\label{used-by}

Connect later dependent concepts here when their explanations use this mechanism.

\subsection{Recall and next depth}\label{recall-and-next-depth}

Reconstruct the example and identify the condition that would invalidate it.

\textbf{Pending:} Target flushing, payload conventions and accumulation instructions require independent checks.

\subsection{Source and navigation}\label{source-and-navigation}

\href{Continuous\%20Notes/Computing\%20Foundations\%20-\%20From\%20Bits\%20to\%20C++\%20and\%20Systems\%20APIs\#Choosing\%20what\%20information\%20to\%20keep}{Continuous teaching story} · \hyperref[appendix-fd090d982a15]{Source coverage ledger} · \href{}{Dependency Map (02-Foundations-II, appendix-4ad47bdd269a)}

\subsection{Range is preserved more than precision}\label{range-is-preserved-more-than-precision}

Bfloat16 keeps the 8-bit exponent width of binary32 but stores only 7 explicit fraction bits. A normal value therefore has 8 bits of significand precision including the hidden leading one. Compared with IEEE binary16, bfloat16 sacrifices substantial local precision to retain a binary32-like exponent range.

That trade is useful for ML tensors whose dynamic range matters more than fine per-value precision. It is poor for algorithms where small relative changes accumulate over many steps. Common mixed-precision kernels store operands in bfloat16 but accumulate products in binary32.

Conversion is not merely truncation if the implementation promises round-to-nearest behavior. Correct rounding must account for discarded bits and tie policy; NaN payloads and infinities need explicit handling. Tests should inspect bit-level round trips, halfway cases, overflow/underflow, NaNs, and error growth in dot products accumulated in bfloat16 versus binary32.
